% (C) 2026 Ralf Stephan, in collaboration with Claude Code.
% Released under CC0 1.0 Universal (public-domain dedication).
% See https://creativecommons.org/publicdomain/zero/1.0/

\documentclass[11pt]{article}
\usepackage[T1]{fontenc}
\usepackage[utf8]{inputenc}
\usepackage{lmodern}
\usepackage[margin=2.5cm]{geometry}
\usepackage{amsmath,amssymb}
\usepackage{booktabs,tabularx}
\usepackage{braket}
\usepackage[hidelinks]{hyperref}

% Lean names and paths: verbatim, breakable at . _ /
\DeclareUrlCommand\lean{\urlstyle{tt}}
\def\UrlBreaks{\do\.\do\_\do\/}
\newcolumntype{L}{>{\raggedright\arraybackslash}X}
\newcolumntype{B}{>{\raggedright\arraybackslash}p{0.42\textwidth}}
\emergencystretch=6em

\title{Lean formalization of an Assumptions of Physics technical brief:\\ log}
\author{}
\date{}

\begin{document}
\maketitle

\noindent
This log records which statements of the technical brief in this directory were checked in Lean,
where the formal version differs from the brief, and any problems found in the brief.
The brief is available at \url{https://assumptionsofphysics.org/resources/briefs}.

\section*{Setup}

\begin{itemize}
\item Lean \texttt{v4.36.0-rc2}; Mathlib at commit \texttt{1b1ccf4b}; physlib (\texttt{master})
  and TauCeti (commit \texttt{e7e7f52}, on the same Mathlib) as further dependencies.
\item Build: \texttt{lake build} in the project root.
\item Every file is checked with \texttt{\#print axioms} on its main results. A result is
  recorded here only if it depends on no axioms beyond \lean{propext},
  \lean{Classical.choice} and \lean{Quot.sound}, and so has no \lean{sorry}.
\item Imports are kept to the Mathlib modules actually needed.
\item Material that belongs in Mathlib is kept in \lean{ForMathlib/}, in Mathlib's directory
  layout (library \lean{ForMathlib}), and imported by the brief files.
\end{itemize}


%======================================================================

\section{Non-uniqueness of the moment problem for Schwartz functions}

\paragraph{Files.}\leavevmode
\begin{itemize}
\item File: \lean{briefs/SchwartzMomentProblem/SchwartzMomentProblem.lean}
\item Library: \lean{SchwartzMomentProblem}
\item Imports: \lean{Mathlib.Analysis.Calculus.BumpFunction.InnerProduct},\newline
  \lean{Mathlib.Analysis.Distribution.SchwartzSpace.Basic},\newline
  \lean{Mathlib.MeasureTheory.Measure.Lebesgue.Basic}
\end{itemize}

\paragraph{Setting.} Schwartz functions $f, g \in \mathcal{S}(\mathbb{R}, \mathbb{C})$
(\lean{SchwartzMap}) with compact supports and $\operatorname{tsupp} f \cap \operatorname{tsupp} g =
\emptyset$, and $\psi_\theta = f + e^{i\theta} g$. The moments are
\[
  \braket{\psi|X^m P^n|\psi} = \int_{\mathbb{R}} \overline{\psi(x)}\, x^m\,
  (-i\hbar)^n\, \psi^{(n)}(x)\, dx,
\]
with $\hbar$ an arbitrary real parameter.

\paragraph{Results.}
\begin{center}
\begin{tabularx}{\textwidth}{@{}BL@{}}
\toprule
Brief & Lean \\
\midrule
cross terms vanish, pointwise & \lean{conj_mul_iteratedDeriv_eq_zero}, \lean{density_psi} \\
eq.~(3): $\braket{\psi_\theta|X^mP^n|\psi_\theta} = \braket{f|X^mP^n|f} + \braket{g|X^mP^n|g}$ & \lean{moment_psi} \\
each $\theta$ gives a different state & \lean{linearIndependent_psi}: $\psi_\theta, \psi_{\theta'}$ linearly independent when $e^{i\theta} \ne e^{i\theta'}$ \\
the bump functions of Fig.~1 & \lean{bump}, \lean{tsupport_bump_subset}, \lean{bump_ne_zero} \\
Theorem~1 & \lean{exists_same_moments}, with $\psi_1 = \psi_0$, $\psi_2 = \psi_\pi$ for bumps centred at $0$ and $3$ \\
\bottomrule
\end{tabularx}
\end{center}

\paragraph{Differences from the brief.}
\begin{itemize}
\item The moment is defined as an integral for every Schwartz function, but the identity
  for $\psi_\theta$ is proved only under the brief's hypothesis that $f$ and $g$ have compact
  support. This makes every integrand continuous with compact support, hence integrable.
\item ``Each $\theta$ defines a different physical state'' is formalized as linear
  independence of $\psi_\theta$ and $\psi_{\theta'}$, which holds exactly when
  $e^{i\theta} \ne e^{i\theta'}$, i.e.\ for $\theta \not\equiv \theta' \pmod{2\pi}$.
\item The closing remark on Schwartz seminorms and the conclusion about local operators
  are not formalized.
\end{itemize}

\paragraph{Problems found in the brief.} No mathematical problems. The argument needs $f$ and
$g$ to be nonzero, which the brief leaves implicit, and ``different state for each $\theta$''
holds only modulo $2\pi$.

\end{document}
